% =============================================================================
% §16b — Extended Timeline Syntax
% Authored: Leslie (Spec Architect), 2026-06-15
% =============================================================================

\section{Extended Timeline Syntax}
\label{sec:extended-timeline}

Mermaid's \texttt{timeline} keyword is intentionally minimal: a sequence of
\texttt{section} groupings each containing \texttt{period : event} lines.
That simplicity makes it portable but leaves the full power of the Timeline
Domain IR~(Section~\ref{sec:timeline-grammar}) unreachable from the text
surface.

The \emph{Extended Timeline} is a strict superset of that syntax.  It uses
the same \texttt{timeline} header keyword and accepts every legal Mermaid
timeline file unchanged; files that use any extended construct are
superset-only and carry a compiler warning (see
Section~\ref{sec:et-degradation}).  The design principle is \emph{one IR,
many diagrams}: a single extended-timeline source compiles to the same
\texttt{IRDocument} regardless of which of the six layout families the author
selects (Section~\ref{sec:et-layouts}), and it inherits all seven contract
themes (Section~\ref{sec:themes}) without modification.

% ---------------------------------------------------------------------------
\subsection{Two-Tier Portability Model}
\label{sec:et-tiers}
% ---------------------------------------------------------------------------

\begin{description}
  \item[Tier 1 — Mermaid-faithful \texttt{timeline}] (unchanged, portable).
        Uses only \texttt{section} and \texttt{period : event} syntax.
        Any valid Mermaid timeline is valid here and renders identically in
        vanilla Mermaid.  The compiler maps these files to the Timeline IR
        via the existing \texttt{parseTimeline} path and defaults the layout
        to \texttt{timeline-columns} to match Mermaid's visual presentation.

  \item[Tier 2 — Extended Timeline] (superset-only, not renderable in
        vanilla Mermaid).  Uses any construct defined in this section: named
        tracks, duration activities, \texttt{@}-directives, sections with
        explicit time ranges, annotations, legend, or axis breaks.  The
        compiler emits a diagnostic warning and proceeds normally; the same
        \texttt{IRDocument} is the output regardless of which layout is
        chosen.
\end{description}

The two tiers share \emph{exactly} the same IR target.  There is no new IR
introduced by the Extended Timeline --- every construct maps to an existing
\texttt{IRDocument} field defined in \texttt{packages/core/src/types.ts} and
validated by \texttt{packages/core/src/schema.ts}.

% ---------------------------------------------------------------------------
\subsection{Frontmatter Configuration}
\label{sec:et-frontmatter}
% ---------------------------------------------------------------------------

The YAML frontmatter block (\texttt{---}\ldots\texttt{---}) configures the
document at parse time.  All keys are optional; unknown keys are warned and
ignored.  The full config surface (Section~\ref{sec:superset-extensions},
\S17.1) is available; the timeline-relevant subset is:

\begin{lstlisting}[language=yaml,caption={Extended timeline frontmatter}]
---
title: "Platform Roadmap 2026"
subtitle: "Engineering org — Q1..Q4"
theme: executive          # any of 7 contract themes
layout: roadmap           # horizontal | vertical-spine | serpentine |
                          # roadmap | gantt | timeline-columns
density: compact          # compact | normal | comfortable
spineSpacing: even        # time | even  (vertical-spine only)
---
\end{lstlisting}

\begin{table}[h]
\centering\small
\caption{Frontmatter keys and their IR targets}
\label{tab:et-frontmatter}
\begin{tabular}{lll}
\toprule
\textbf{Key} & \textbf{IR field} & \textbf{Notes} \\
\midrule
\texttt{title}       & \texttt{metadata.title}        & Already supported in Tier 1 \\
\texttt{subtitle}    & \texttt{metadata.subtitle}     & Already supported in Tier 1 \\
\texttt{theme}       & \texttt{metadata.theme}        & 7 contract theme names valid \\
\texttt{layout}      & \texttt{metadata.layout}       & 6 values; see Table~\ref{tab:et-layouts} \\
\texttt{density}     & Theme-resolution config        & Not an IR \texttt{Metadata} field; \\
                     &                                & passed to \texttt{resolveContractTheme} \\
\texttt{spineSpacing}& \texttt{RenderOptions.spineSpacing} & \texttt{time} (proportional) or \\
                     &                                & \texttt{even} (infographic); \\
                     &                                & only affects \texttt{vertical-spine} \\
\bottomrule
\end{tabular}
\end{table}

\noindent
\textbf{Note on \texttt{layout} schema gap.}  As of 2026-06-15,
\texttt{metadata.layout} in \texttt{schema.ts} accepts only four values
(\texttt{horizontal}, \texttt{vertical-spine}, \texttt{serpentine},
\texttt{roadmap}); the TypeScript type in \texttt{types.ts} correctly
includes all six.  The Zod schema must be extended to add \texttt{gantt} and
\texttt{timeline-columns} before those values can be round-tripped through
JSON Schema validation.  This is a known schema gap, not a design gap.

% ---------------------------------------------------------------------------
\subsection{Grammar Sketch}
\label{sec:et-grammar}
% ---------------------------------------------------------------------------

The following EBNF-style grammar covers the Extended Timeline body (after
frontmatter extraction).  Whitespace between tokens is insignificant except
that indented lines nest under the nearest enclosing construct.  Two-space
indentation is conventional; the parser accepts any consistent depth $\ge
1$.

\begin{lstlisting}[caption={Extended Timeline EBNF sketch}]
extended-timeline ::=
    'timeline' [ 'extended' ] EOL
    stmt*

stmt ::=
      title-stmt
    | track-stmt
    | section-stmt
    | annotation-stmt
    | legend-stmt
    | break-stmt

(* ---- top-level declarations ---- *)
title-stmt    ::= 'title' text EOL

track-stmt    ::= 'track' track-name EOL
                  INDENT track-item*

track-item    ::= activity-stmt | milestone-stmt

(* ---- activities ---- *)
activity-stmt ::= label ':' date-expr EOL
                  ( INDENT directive )*

date-expr     ::= irdate '..' irdate   (* duration form *)
               |  irdate               (* single-date (span) form *)

(* ---- milestones ---- *)
milestone-stmt ::= '@milestone' label ':' irdate EOL
                   ( INDENT directive )*

(* ---- @-directives (sub-lines under activity or milestone) ---- *)
directive     ::= '@status'      status-value EOL
               |  '@progress'    integer      EOL   (* 0-100 *)
               |  '@shape'       shape-token  EOL
               |  '@icon'        icon-name    EOL
               |  '@color'       css-color    EOL
               |  '@description' text         EOL

(* ---- sections ---- *)
section-stmt  ::= 'section' label
                  ( '[' irdate '..' irdate ']' )? EOL

(* ---- annotations ---- *)
annotation-stmt ::= 'annotation' DQUOTED-TEXT
                    '@' annotation-target EOL

annotation-target ::= irdate '..' irdate    (* period annotation *)
                    | irdate                (* point annotation  *)

(* ---- legend ---- *)
legend-stmt   ::= 'legend' ( 'show' | 'hide' )? EOL

(* ---- axis breaks ---- *)
break-stmt    ::= 'break' irdate '..' irdate EOL

(* ---- primitives ---- *)
irdate   ::= (* IR date string per §21 §4 date model;
               examples: 2026-Q2  2026-06-01  FY26-Q3  now  tbd *)

status-value ::=
    'planned' | 'in-progress' | 'done' |
    'at-risk'  | 'blocked'    | 'cancelled' | 'tentative'

shape-token  ::= (* theme-resolved icon name, e.g. 'diamond' 'star' 'flag' *)

label        ::= (* non-empty text, no leading colon *)
track-name   ::= (* non-empty text *)
\end{lstlisting}

\paragraph{Keyword \texttt{extended}.}  The optional \texttt{extended}
modifier on the opening line is a forward hint.  When present, the compiler
suppresses the ``extended features detected'' warning (the author has
explicitly opted in).  When absent, the warning is emitted on the first use
of any Tier-2 construct.

\paragraph{Indentation and nesting.}
A \texttt{track} header owns all activity and milestone lines at the next
indentation level until a dedent or another \texttt{track}/\texttt{section}/
\texttt{annotation} at the same level appears.  \texttt{@}-directives attach
to the \emph{immediately preceding} activity or milestone and must be
indented one level deeper than that activity.

% ---------------------------------------------------------------------------
\subsection{Construct-to-IR Mapping}
\label{sec:et-ir-mapping}
% ---------------------------------------------------------------------------

Every Extended Timeline construct maps to an existing \texttt{IRDocument}
field.  Table~\ref{tab:et-ir-mapping} gives the authoritative mapping using
the canonical TypeScript field names from \texttt{types.ts} in
\texttt{packages/core/src/}.

\begin{table}[h]
\centering\small
\caption{Extended Timeline syntax $\rightarrow$ IRDocument field mapping}
\label{tab:et-ir-mapping}
\begin{tabular}{p{3.8cm}p{4.2cm}p{5.0cm}}
\toprule
\textbf{DSL construct} & \textbf{IR field} & \textbf{Notes} \\
\midrule
\texttt{track Name}
  & \texttt{Track \{ id, label \}}
  & \texttt{id} = sanitized kebab-case slug of Name \\
\addlinespace
\texttt{Label: start..end}
  & \texttt{Activity \{ track, label,} \texttt{start, end \}}
  & Both \texttt{start} and \texttt{end} are \texttt{IRDate} strings \\
\addlinespace
\texttt{Label: date}
  & \texttt{Activity \{ track, label, span \}}
  & Single-date form; \texttt{Activity.span} is the period identifier \\
\addlinespace
\texttt{@status value}
  & \texttt{Activity.status} or \texttt{Milestone.status}
  & Must be one of the 7 \texttt{Status} enum values \\
\addlinespace
\texttt{@progress N}
  & \texttt{Activity.progress}
  & Author writes integer 0--100; parser stores $N \div 100 \in [0,1]$ \\
\addlinespace
\texttt{@milestone Label: date}
  & \texttt{Milestone \{ id, label, date, track \}}
  & \texttt{track} is the enclosing \texttt{track} block's ID \\
\addlinespace
\texttt{@shape token}
  & \texttt{Milestone.icon}
  & No \texttt{shape} field in IR; \texttt{icon} is the closest proxy.
    Exact marker geometry is theme-resolved.
    \emph{(Gap: see \S\ref{sec:et-gaps})} \\
\addlinespace
\texttt{section Name [s..e]}
  & \texttt{Section \{ id, label, time\_range \}}
  & \texttt{time\_range.start} and \texttt{time\_range.end} from bracket \\
\addlinespace
\texttt{annotation "text" @ date}
  & \texttt{Annotation \{ type:'callout', text, date \}}
  & Point annotation; \texttt{type} = \texttt{'callout'} \\
\addlinespace
\texttt{annotation "text" @ s..e}
  & \texttt{Annotation \{ type:'period', text, start, end \}}
  & Range annotation; \texttt{type} = \texttt{'period'} \\
\addlinespace
\texttt{legend show}
  & \texttt{Legend \{ show: true \}}
  & Renderer auto-builds entries from \texttt{Status}/\texttt{category} data \\
\addlinespace
\texttt{break s..e}
  & \texttt{Metadata.axis\_breaks: [\{ from, to \}]}
  & Collapses dead calendar gaps to a \texttt{//} notch \\
\bottomrule
\end{tabular}
\end{table}

\paragraph{The \texttt{Status} enum.}
The seven legal \texttt{@status} values (from \texttt{types.ts}) are:
\texttt{planned}, \texttt{in-progress}, \texttt{done}, \texttt{at-risk},
\texttt{blocked}, \texttt{cancelled}, \texttt{tentative}.
The theme contract maps each to a role-palette color
(Section~\ref{sec:contract-palette}).

\paragraph{The \texttt{IRDate} format.}
Both \texttt{start..end} operands and standalone \texttt{date} values must
conform to the date model defined in Section~\ref{sec:date-model} and the
regex in \texttt{schema.ts}.  Legal examples: \texttt{2026-06-01},
\texttt{2026-Q2}, \texttt{2026-H1}, \texttt{FY26-Q3}, \texttt{now},
\texttt{tbd}, \texttt{\~{}2027-Q1}.

\paragraph{\texttt{@progress} integer-to-float mapping.}
The IR stores \texttt{Activity.progress} as a decimal in $[0, 1]$ (type
\texttt{number[0..1]}).  The DSL accepts an author-friendly integer percent:
$\texttt{@progress}\ N \mapsto \texttt{progress} = N / 100$.  The parser
clamps the value to $[0, 1]$; values outside $[0, 100]$ emit a warning.

% ---------------------------------------------------------------------------
\subsection{One IR, Many Layouts}
\label{sec:et-layouts}
% ---------------------------------------------------------------------------

The canonical claim of the Extended Timeline is the \emph{reach principle}
(Section~\ref{sec:themes-reach}): one \texttt{IRDocument} can be rendered
under any of the six layout families without changing the source.  The
author sets \texttt{layout:} in the frontmatter; switching layouts requires
editing \emph{only that line}.

\begin{table}[h]
\centering\small
\caption{The six layout families and their extended-timeline affordances}
\label{tab:et-layouts}
\begin{tabular}{lp{4.2cm}p{6.5cm}}
\toprule
\textbf{\texttt{layout} value} & \textbf{Character} & \textbf{Best-fit use case} \\
\midrule
\texttt{horizontal}       & Time axis left$\rightarrow$right, track rows
  & Dense multi-track engineering schedules; high information density \\
\texttt{vertical-spine}   & Central spine, L/R alternating entries
  & Narrative timelines; investor decks; few milestones \\
\texttt{serpentine}       & Boustrophedon winding path
  & Journey maps; annual reviews with a clear chronological story \\
\texttt{roadmap}          & Band-per-track, month/quarter columns
  & Product roadmaps; portfolio views; the \emph{executive} sweet spot \\
\texttt{gantt}            & Mermaid-faithful Gantt: task bars, date axis
  & Project management; dependency-heavy schedules \\
\texttt{timeline-columns} & Section-colored header bands, evenly-spaced columns
  & Mermaid-compatible presentation; events stacked under periods \\
\bottomrule
\end{tabular}
\end{table}

\paragraph{The $7 \times 6 = 42$ look matrix.}
Any of the seven contract themes (\texttt{executive}, \texttt{midnight},
\texttt{blueprint}, \texttt{editorial}, \texttt{terminal}, \texttt{pastel},
\texttt{mono}) combined with any of the six layouts produces a distinct
visual presentation from a single source file.  Themes drive typography,
palette, and connector style (Section~\ref{sec:themes-drives-layout});
layouts drive geometry and axis orientation
(Section~\ref{sec:timeline-rendering}).  This \emph{combination matrix} is
the primary differentiator of the Extended Timeline over Mermaid's fixed
rendering.

\paragraph{The \texttt{poster} bridge.}
A \texttt{poster} (Section~\ref{sec:superset-extensions}) can reference the
same extended-timeline source under multiple layouts in adjacent cells,
producing a side-by-side showcase of the $42$-look matrix in a single
document.

\paragraph{Dimension guard.}
The layout engine enforces a dimension guard (height $\le 5000$\,px; aspect
ratio $\le 4:1$; see \texttt{gallery-dimensions.test.ts}).  Pathological
combinations --- for example, \texttt{vertical-spine} over a multi-decade
span with default \texttt{spineSpacing:time} --- trigger a render-time
warning and a concrete suggestion:

\begin{quote}
\textit{``vertical-spine with a time-proportional spine exceeds safe render
dimensions.  Consider \texttt{spineSpacing: even} or switching to
\texttt{layout: roadmap} / \texttt{timeline-columns} for multi-decade
spans.''}
\end{quote}

This is implemented in \texttt{packages/core/src/grammars/timeline/\allowbreak
vertical-spine.ts} and was validated against the full gallery in the
dimension-guard decision (2026-06-15).

% ---------------------------------------------------------------------------
\subsection{Degradation and Portability Contract}
\label{sec:et-degradation}
% ---------------------------------------------------------------------------

\begin{quote}
\textbf{Extended Timeline Portability Contract (2026-06-15).}

\begin{enumerate}
  \item A file using only Tier-1 constructs (\texttt{section} +
        \texttt{period : event}) is byte-compatible with vanilla Mermaid and
        produces no warnings from this compiler.
  \item A file using any Tier-2 construct is superset-only.  The compiler
        emits a \texttt{WARNING} at parse time:
        \textit{``Extended timeline features detected --- output is not
        renderable by vanilla Mermaid.  Use Tier-1 syntax for portable
        files.''} and continues compilation normally.
  \item The warning is \emph{suppressed} when the opening line reads
        \texttt{timeline extended}, signalling explicit opt-in.
  \item All Tier-2 constructs map to existing \texttt{IRDocument} fields.
        No new IR is introduced.  Switching from a Tier-2 file back to
        Tier-1 syntax produces a valid (reduced) document with no schema
        changes.
  \item Vanilla Mermaid renderers encountering a Tier-2 file will fail to
        parse it predictably (they do not silently corrupt output); the
        failure is a clear signal to switch to this compiler.
\end{enumerate}
\end{quote}

% ---------------------------------------------------------------------------
\subsection{Worked Examples}
\label{sec:et-examples}
% ---------------------------------------------------------------------------

\subsubsection{Full Extended Timeline}

The following example exercises every Tier-2 construct and shows the
frontmatter configuration.

\begin{lstlisting}[caption={Full extended timeline source}]
---
title: "Platform Engineering Roadmap"
subtitle: "2026 — all tracks"
theme: executive
layout: roadmap
density: normal
---
timeline extended

section Infrastructure [2026-Q1 .. 2026-Q2]
section Product [2026-Q2 .. 2026-Q4]

track Platform
  Kubernetes Upgrade: 2026-Q1 .. 2026-Q2
    @status done
    @progress 100
    @description "Migrate all clusters to k8s 1.32"
  Service Mesh Rollout: 2026-Q2 .. 2026-Q3
    @status in-progress
    @progress 40
  @milestone Hard Freeze: 2026-06-30
    @status planned
    @shape diamond

track API
  GraphQL Gateway: 2026-Q1 .. 2026-Q3
    @status in-progress
    @progress 60
  @milestone v3 Launch: 2026-09-01
    @status planned
    @shape star

track Mobile
  iOS Redesign: 2026-Q2
    @status planned
  Android Parity: 2026-Q3 .. 2026-Q4
    @status tentative

annotation "PCI-DSS audit window" @ 2026-06-01 .. 2026-06-30
annotation "Board review" @ 2026-Q3

break 2026-12-20 .. 2027-01-05

legend show
\end{lstlisting}

\subsubsection{Same Data, Multiple Layouts}

Switching \texttt{layout:} in the frontmatter (one line) re-renders the
identical \texttt{IRDocument} under a different presentation family.

\begin{lstlisting}[caption={Same source, four layout variants}]
# Layout 1: Gantt — dependency-aware task bars
layout: gantt

# Layout 2: Roadmap — band-per-track, quarter columns
layout: roadmap

# Layout 3: Serpentine — boustrophedon journey narrative
layout: serpentine

# Layout 4: Timeline-columns — Mermaid-compatible column bands
layout: timeline-columns
\end{lstlisting}

\noindent
Combined with seven contract themes, this yields 42 distinct visual
presentations from one source.  A \texttt{poster} cell grid can render
selected layout-theme combinations side by side:

\begin{lstlisting}[language=yaml,caption={Poster showcasing two layout variants of the same source}]
---
theme: executive
layout: grid 1x2
---
poster "Roadmap two ways"
  cell [0,0]: timeline extended
    # ... (roadmap layout via frontmatter override)
  cell [0,1]: timeline extended
    # ... (gantt layout via frontmatter override)
\end{lstlisting}

% ---------------------------------------------------------------------------
\subsection{Known IR Gaps and Deferred Questions}
\label{sec:et-gaps}
% ---------------------------------------------------------------------------

The following items are gaps between what the Extended Timeline syntax
expresses and what the current IR (2026-06-15) can cleanly carry.  They are
flagged here for future IR revision, not as blockers for the syntax
definition.

\begin{description}
  \item[\texttt{@shape} has no IR field.]
        \texttt{Milestone} in \texttt{types.ts} carries \texttt{icon?:
        string} and \texttt{category?: string} but no \texttt{shape} field.
        The parser maps \texttt{@shape diamond} to \texttt{Milestone.icon =
        'diamond'}, conflating pictogram identity with marker geometry.
        Future IR revision should add \texttt{Milestone.shape?: enum\{diamond,
        circle, square, star, flag\}} and separate it from \texttt{icon}.

  \item[\texttt{gantt} and \texttt{timeline-columns} missing from Zod schema.]
        \texttt{metadata.layout} in \texttt{schema.ts} lists only four values
        (\texttt{horizontal}, \texttt{vertical-spine}, \texttt{serpentine},
        \texttt{roadmap}).  The TypeScript type in \texttt{types.ts} correctly
        includes all six.  The schema must be extended to add \texttt{gantt}
        and \texttt{timeline-columns} for these layouts to survive
        JSON Schema round-trip validation.

  \item[\texttt{density} is not an \texttt{IRDocument} field.]
        Density is resolved at theme time via \texttt{resolveContractTheme}
        and never written into the IR.  A document authored at
        \texttt{density: compact} is indistinguishable from one at
        \texttt{density: normal} after IR round-trip.  Whether density should
        be promoted to \texttt{Metadata} is an open question for the Tier-2
        contract evolution (Section~\ref{sec:themes-contract}).

  \item[\texttt{legend} auto-entry generation is unspecified.]
        The IR supports \texttt{Legend.entries} (explicit \texttt{LegendEntry}
        objects), but the parser-level behaviour of \texttt{legend show}
        without explicit entries (auto-build from status/category data) is a
        rendering convention, not an IR invariant.  The spec defers the
        auto-entry algorithm to the layout engine documentation
        (Section~\ref{sec:timeline-rendering}).

  \item[\texttt{@milestone} without an enclosing \texttt{track}.]
        If \texttt{@milestone} appears outside any \texttt{track} block, the
        parser should assign \texttt{Milestone.track = undefined} (the IR
        permits this).  The layout engine then renders the milestone on the
        global axis rather than a swimlane.  This behaviour should be
        explicitly validated.
\end{description}
